\chapter{Benchmark-Design}
\label{ch:benchmark-design}

Dieses Kapitel beschreibt f{\"u}r jeden der f{\"u}nf Benchmarks im Detail, was genau gemessen wird, wie der Messpunkt technisch implementiert ist, welche Fehlerquellen dabei kontrolliert werden m{\"u}ssen und warum die jeweilige Metrik f{\"u}r die Forschungsfrage relevant ist. Vor den DWT-basierten Messungen wird \lstinline[style=fileinline]|k_busy_wait(100)| gemessen. Bei 64 MHz liegt der nominelle Wert bei 6400 CPU-Zyklen inklusive des Eigenaufwands der Funktion. Als Abnahmekriterium gelten \lstinline[style=fileinline]!|mean - 6400| < 1 %! und eine Standardabweichung unter zwei Zyklen. Der Test erkennt dabei grobe Fehler von Taktquelle, DWT oder Abzug des Messaufwands, beweist jedoch nicht die exakte Lage asynchroner Messpunkte oder die Abwesenheit von Interrupts in anderen Benchmarks.

\section{Kontextwechsel-Latenz}
\begin{lstlisting}[style=basecode,language=C]
/* Variante 1: sem_wakeup */
...
timing_t t0 = bench_now();
k_sem_give(&wake_sem);
/* Worker hat t_stamp gesetzt und blockiert wieder */
bench_stats_add(&st, bench_elapsed(t0, t_stamp) - bench_overhead_cycles());
...
/* Variante 2: yield (im Worker-Thread gemessen) */
...
timing_t t1 = bench_now();
bench_stats_add(&yield_stats, bench_elapsed(t_stamp, t1) - bench_overhead_cycles());
k_yield();
\end{lstlisting}
\textbf{Definition} Gemessen wird die Zeit zwischen dem Zeitstempel direkt vor der ausl{\"o}senden Kerneloperation und dem fr{\"u}h im Zielthread erfassten Zeitstempel. Das Intervall umfasst dabei die bis zu diesem Messpunkt durchlaufenen Verwaltungs- und Umschaltoperationen. Diese Zeit umfasst auf einem Cortex-M die Verwaltung des ausl{\"o}senden Kernelaufrufs, die Scheduler-Entscheidung und das Sichern und Laden des Thread-Kontexts. Laut Zephyr-Dokumentation ist ein Kontextwechsel unter anderem genau dann m{\"o}glich, wenn ein Thread {\"u}ber \lstinline[style=fileinline]|k_sem_take()| in den wartenden Zustand {\"u}bergeht oder {\"u}ber \lstinline[style=fileinline]|k_sem_give()|/ \lstinline[style=fileinline]|k_yield()| einen anderen Thread bereit macht bzw. die CPU freiwillig abgibt. Diese Reschedule Points bilden das Grundkonzept der beiden gemessenen Varianten.

\textbf{Messmethode} Diese Untersuchung unterscheidet zwischen zwei Varianten des Kontextwechsels:

\begin{longtable}{p{0.235\linewidth}p{0.235\linewidth}p{0.235\linewidth}p{0.235\linewidth}}
\toprule
Variante & Aufbau & Messpunkt Start & Messpunkt Ende \\
\midrule
\lstinline[style=fileinline]|ctx_switch_sem_wakeup| & Ein h{\"o}herpriorer Worker-Thread blockiert auf \lstinline[style=fileinline]|k_sem_take()|, der niederpriore Haupt-Thread ruft k\_sem\_give() auf & Direkt vor \lstinline[style=fileinline]|k_sem_give()| & Zeitstempel als erste C-seitige Aktion nach R{\"u}ckkehr aus \lstinline[style=fileinline]|k_sem_take()| \\
\lstinline[style=fileinline]|ctx_switch_yield| & Zwei gleichpriore Threads {\"u}bergeben sich kooperativ {\"u}ber \lstinline[style=fileinline]|k_yield()| & Direkt vor \lstinline[style=fileinline]|k_yield()| & Zeitstempel am Beginn der entsprechenden Worker-Iteration \\
\bottomrule
\end{longtable}

Bei der ersten Variante garantiert die h{\"o}here Priorit{\"a}t des Workers, dass dieser nach \lstinline[style=fileinline]|k_sem_give()| sofort l{\"a}uft, bevor der Haupt-Thread weiterl{\"a}uft. Laut Zephyr-Dokumentation wird ein Semaphor an den wartenden Thread mit der h{\"o}chsten Priorit{\"a}t vergeben. Bei der zweiten Variante wechseln sich beide Threads ab, da die Zeitscheiben-Rotation in der Minimal-Konfiguration deaktiviert ist (siehe 4.4.1). Dabei reiht \lstinline[style=fileinline]|k_yield()| den aufrufenden Thread laut Zephyr-Dokumentation ans Ende der Liste der bereiten Threads mit gleicher Priorit{\"a}t ein und {\"u}bergibt die CPU dann sofort an den n{\"a}chsten. Beide Varianten werden hierbei auch von Zephyrs eigenem Referenz-Benchmark \lstinline[style=fileinline]|tests/benchmarks/latency_measure| erfasst, was einen externen Validierungsvergleich erm{\"o}glicht.

\textbf{Fehlerquellen} Interrupts k{\"o}nnen die gemessene Zeit verl{\"a}ngern. Tickless-Betrieb verringert periodische Systemtimer-Unterbrechungen, schlie{\ss}t andere Interruptquellen jedoch nicht aus. Eine geringe Selbsttest-Streuung belegt deshalb nicht, dass w{\"a}hrend den Benchmark-Reihen keine Interrupts auftreten. Die erste Worker-Ankunft l{\"a}uft dabei {\"u}ber den Thread-Startpfad und wird durch die Warm-up-Durchl{\"a}ufe verworfen. MPU und hardwaregest{\"u}tzter Stack-Schutz sind in allen drei Benchmark-Profilen deaktiviert. Deren Zusatzkosten werden ohne separaten Vergleichslauf nicht quantifiziert.

\textbf{Relevanz} Kontextwechsel wurden als Metrik genommen, weil sie die Grundkosten jeder Multithreading-Architektur sind und bestimmen, wie schnell ein RTOS auf Ereignisse reagieren kann. Zudem bestimmen sie, wie fein die Aufgaben eines Systems in einzelne Threads zerlegt werden k{\"o}nnen, ohne dass der Wechsel selbst zum limitierenden Faktor wird.

\section{Interrupt-Latenz}
\begin{lstlisting}[style=basecode,language=C]
static void trigger_isr(const void *arg)
{
    ...
    isr_stamp = bench_now();
    isr_count++;
}
...
uint32_t expected = isr_count + 1;
timing_t t0 = bench_now();
NVIC_SetPendingIRQ(TRIGGER_IRQN);
while (isr_count != expected) { /* aktives Warten */ }
bench_stats_add(&st, bench_elapsed(t0, isr_stamp) - bench_overhead_cycles());
\end{lstlisting}
\textbf{Definition} Hierbei handelt es sich um die Zeit zwischen dem softwareseitigen Ausl{\"o}sen eines Interrupts und dem Zeitstempel am Anfang des registrierten C-Handlers. Der Exception-Eintritt, der Zephyr-Interrupt-Wrapper und der Compilerprolog liegen vor diesem Messpunkt und sind daher Teil des Messwerts. Eine ISR ist laut Zephyr-Dokumentation eine Funktion, die asynchron als Reaktion auf einem Interrupt ausgef{\"u}hrt wird und dabei die Ausf{\"u}hrung des aktuellen Threads unterbricht. Ziel ist es hierbei, mit m{\"o}glichst geringem Messaufwand reagieren zu k{\"o}nnen.

\textbf{Messmethode} F{\"u}r einen softwareseitig kontrollierten Ausl{\"o}sezeitpunkt benutzt der Benchmark die IRQ-Leitung von \lstinline[style=fileinline]|TIMER1| und setzt deren Pending-Bit {\"u}ber \lstinline[style=fileinline]|NVIC_SetPendingIRQ()|. Die \lstinline[style=fileinline]|TIMER1|-Peripherie muss daf{\"u}r nicht selbst als laufender Zeitgeber betrieben werden. Voraussetzung ist hierbei, dass die ausgew{\"a}hlte IRQ-Leitung im untersuchten Build nicht wo anders benutzt wird. Der Endzeitstempel wird am Anfang des registrierten C-Handlers erfasst. Exception-Eintritt, Zephyr-Interrupt-Wrapper und Compiler-Prolog liegen vor diesem Messpunkt und sind Teil des Messwerts.

\textbf{Fehlerquellen} Der Trigger-IRQ besitzt die numerische Priorit{\"a}t 3 und kann Thread-Mode-Code unterbrechen. Er kann durch Interrupts h{\"o}herer Dringlichkeit unterbrochen oder durch Interruptsperren verz{\"o}gert werden. Die Warteschleife wird erst nach dem Endzeitstempel ausgef{\"u}hrt und geht deshalb nicht in die Zeitdifferenz ein. Die Wahl von \lstinline[style=fileinline]|TIMER1| und seiner IRQ-Nummer ist plattformspezifisch. Dabei geh{\"o}rt \lstinline[style=fileinline]|NVIC_SetPendingIRQ()| zur CMSIS-Core-Schnittstelle. Der dazugeh{\"o}rige Typ hei{\ss}t \lstinline[style=fileinline]|IRQn_Type|, und der benutzte Header ist \lstinline[style=fileinline]|<nrfx.h>|.

\textbf{Relevanz }Ein RTOS verspricht, im Gegensatz zu einem klassischen Betriebssystem, Determinismus. Die Interrupt-Latenz zeigt, wie stark Kernel-Funktionen, wie bspw. das Logging-Subsystem, den zeitkritischen Reaktionspfad eines Systems beeinflussen.

\section{Timer-Jitter}
\begin{lstlisting}[style=basecode,language=C]
static void timer_cb(struct k_timer *t)
{
    uint32_t now_cycles = sys_clock_cycle_get_32();

    if (have_prev) {
        uint32_t delta = now_cycles - prev_cycles;

        if (count >= CONFIG_BENCHLIB_WARMUP) {
            bench_stats_add(&st, delta);
        }
        count++;
    } else {
        have_prev = true;
    }
    prev_cycles = now_cycles;

    if (count >= CONFIG_BENCHLIB_WARMUP + CONFIG_BENCHLIB_ITERATIONS) {
        k_timer_stop(t);
        done = true;
    }
}
\end{lstlisting}
\textbf{Definition:} Erfasst wird die beobachtete Periode eines periodischen Timers. Jitter wird anhand der Abweichungen und Streuung dieser Perioden im Vergleich zu, der in \lstinline[style=fileinline]|TIMERMETA| angegebenen nominellen RTC-Periode beurteilt.

\textbf{Messmethode:} Der \lstinline[style=fileinline]|k_timer| besitzt eine angeforderte Periode von 10 ms. Im Callback wird \lstinline[style=fileinline]|sys_clock_cycle_get_32()| gelesen. Das Delta von zwei aufeinanderfolgenden Reads ist der Abstand der Messpunkte in RTC-Systemtimerzyklen. Der Messpunkt liegt erst im C-Callback und beinhaltet deshalb auch Verz{\"o}gerungen vor seiner Ausf{\"u}hrung.

\textbf{Fehlerquellen:} Der erste Callback liefert noch keinen Messwert, da keine Vorg{\"a}nger existieren. Danach werden zudem zehn Perioden verworfen. Die Aufl{\"o}sung betr{\"a}gt etwa 30,52 \textmu{}s. Interrupts mit h{\"o}herer Dringlichkeit, Interruptsperren und andere Kernelaktivit{\"a}t k{\"o}nnen den Callback verz{\"o}gern. Ein aktiviertes Logging-Subsystem ohne erzeugte Log-Nachrichten erzeugt dabei nicht automatisch fortlaufende UART-Last.

\textbf{Relevanz:} Die Messung zeigt, wie stabil periodische Callback-Messpunkte unter dem jeweiligen kompletten Konfigurationsprofil sind.

\section{Heap-Allokationsgeschwindigkeit}
\begin{lstlisting}[style=basecode,language=C]
...
bench_stats_reset(&alloc_st);
    bench_stats_reset(&free_st);

    for (int i = 0; i < CONFIG_BENCHLIB_WARMUP + CONFIG_BENCHLIB_ITERATIONS; i++) {
        timing_t t0 = bench_now();
        void *p = k_malloc(ALLOC_SIZE);
        timing_t t1 = bench_now();

        timing_t t2 = bench_now();
        k_free(p);
        timing_t t3 = bench_now();

        if (i >= CONFIG_BENCHLIB_WARMUP) {
            bench_stats_add(&alloc_st, bench_elapsed(t0, t1) -
                            bench_overhead_cycles());
            bench_stats_add(&free_st, bench_elapsed(t2, t3) -
                           bench_overhead_cycles());
        }
    }
...
\end{lstlisting}
\textbf{Definition }Bei dieser Metrik wurde die Zeit f{\"u}r die Allokation durch \lstinline[style=fileinline]|k_malloc()| bzw. Freigabe durch \lstinline[style=fileinline]|k_free()| eines Speicherblocks mit fester Gr{\"o}{\ss}e aus dem Zephyr-System-Heap gemessen. Die Gr{\"o}{\ss}e des Blocks wurde hierbei auf 64 Byte festgelegt

\textbf{Messmethode }Pro Iteration wird ein Block angefordert und direkt danach wieder freigegeben, anstatt mehrere Bl{\"o}cke gleichzeitig zu halten. Dabei verhalten sich \lstinline[style=fileinline]|k_malloc()| und \lstinline[style=fileinline]|k_free()| laut Zephyr-Dokumentation semantisch wie die gleichnamigen Standard-C-Funktionen. Sie arbeiten zudem auf einem einzigen, {\"u}ber \lstinline[style=fileinline]|CONFIG_HEAP_MEM_POOL_SIZE| definierten System-Heap. Dabei werden Allokation und Freigabe als zwei getrennte Zeiten gemessen, jeweils direkt vor und nach dem entsprechenden Aufruf.

\textbf{Fehlerquellen }Das direkte Freigeben h{\"a}lt den Systemheap nahezu leer und untersucht damit einen vorteilhaften, unfragmentierten Zustand. Der Fragmentierungsbenchmark untersucht anschlie{\ss}end nur die Allokierbarkeit eines getrennten \lstinline[style=fileinline]|k_heap|. Er misst jedoch nicht die Laufzeit von \lstinline[style=fileinline]|k_malloc()| unter Fragmentierung. Aussagen {\"u}ber Worst-Case-Allokationszeiten oder den Einfluss eines fragmentierten Systemheaps lassen sich aus diesen beiden Benchmarks deshalb nicht treffen.

\textbf{Relevanz }Der Allokationsbenchmark pr{\"u}ft den R{\"u}ckgabewert von \lstinline[style=fileinline]|k_malloc()| nicht. Seine Werte d{\"u}rfen nur als erfolgreiche 64-Byte-Allokationen interpretiert werden, wenn der Erfolg im Messcode zudem best{\"a}tigt wird.

\section{Heap-Fragmentierung}
\begin{lstlisting}[style=basecode,language=C]
static void print_stats(const char *label)
{
    printk("HEAPSTATS,%s,requested=%zu,unrequested=%zu\n",
       label, bytes_allocated,
       (size_t)FRAG_HEAP_SIZE - bytes_allocated);
}
/* Bisektion: groessten Block finden, der aktuell am Stueck allokierbar
 * ist. */
static size_t probe_largest_block(void)
{
    size_t lo = 0, hi = FRAG_HEAP_SIZE;

    while (lo + 8 < hi) {
        size_t mid = (lo + hi) / 2;
        void *p = k_heap_alloc(&frag_heap, mid, K_NO_WAIT);

        if (p) {
            k_heap_free(&frag_heap, p);
            lo = mid;
        } else {
            hi = mid;
        }
    }
    return lo;
}
\end{lstlisting}
\textbf{Definition:} Ausgegeben werden erfolgreich angeforderte Nutzbytes, die rechnerische Differenz zur bereitgestellten Heap-Gr{\"o}{\ss}e und die gr{\"o}{\ss}te im 8-Byte-Suchraster erfolgreich gepr{\"u}fte Einzelanforderung. Dabei ist \lstinline[style=fileinline]|unrequested| keine Messung von freien Allokatorbytes, da Metadaten, Ausrichtung und Rundung nicht beachtet werden.

\textbf{Messmethode:} Von 32 angeforderten 128-Byte-Bl{\"o}cken werden nur erfolgreiche Allokationen verbucht. Beim Freigeben werden NULL-Eintr{\"a}ge {\"u}bersprungen. Die Bisektions-Probe allokiert und gibt Testbl{\"o}cke direkt wieder frei. Wegen der Abbruchbedingung liefert sie eine praktische N{\"a}herung im 8-Byte-Raster und nicht den mathematisch genau gr{\"o}{\ss}ten Freiblock.

\textbf{Fehlerquelle:} Die Probe benutzt den untersuchten Allokator selbst. Nach jedem erfolgreichen Versuch sind keine weiteren Nutzbl{\"o}cke mehr angefordert. Daraus folgt aber nicht unbedingt, dass die internen Verwaltungsstrukturen unver{\"a}ndert sind.

\textbf{Relevanz:} Externe Fragmentierung kann dazu f{\"u}hren, dass eine gro{\ss}e zusammenh{\"a}ngende Anforderung nicht erfolgreich ist, obwohl die Summe der nicht angeforderten bzw. an sich verf{\"u}gbaren Bereiche gr{\"o}{\ss}er ist. Dieses Risiko entsteht durch dynamische Allokation und begrenzten zusammenh{\"a}ngenden Speicher, nicht nur durch das Fehlen einer MMU.
